% Options for packages loaded elsewhere
\PassOptionsToPackage{unicode}{hyperref}
\PassOptionsToPackage{hyphens}{url}
\documentclass[
]{ltjsarticle}
\usepackage{xcolor}
\usepackage{amsmath,amssymb}
\setcounter{secnumdepth}{-\maxdimen} % remove section numbering
\usepackage{iftex}
\ifPDFTeX
  \usepackage[T1]{fontenc}
  \usepackage[utf8]{inputenc}
  \usepackage{textcomp} % provide euro and other symbols
\else % if luatex or xetex
  \usepackage{unicode-math} % this also loads fontspec
  \defaultfontfeatures{Scale=MatchLowercase}
  \defaultfontfeatures[\rmfamily]{Ligatures=TeX,Scale=1}
\fi
\usepackage{lmodern}
\ifPDFTeX\else
  % xetex/luatex font selection
\fi
% Use upquote if available, for straight quotes in verbatim environments
\IfFileExists{upquote.sty}{\usepackage{upquote}}{}
\IfFileExists{microtype.sty}{% use microtype if available
  \usepackage[]{microtype}
  \UseMicrotypeSet[protrusion]{basicmath} % disable protrusion for tt fonts
}{}
\makeatletter
\@ifundefined{KOMAClassName}{% if non-KOMA class
  \IfFileExists{parskip.sty}{%
    \usepackage{parskip}
  }{% else
    \setlength{\parindent}{0pt}
    \setlength{\parskip}{6pt plus 2pt minus 1pt}}
}{% if KOMA class
  \KOMAoptions{parskip=half}}
\makeatother
\usepackage{longtable,booktabs,array}
\newcounter{none} % for unnumbered tables
\usepackage{calc} % for calculating minipage widths
% Correct order of tables after \paragraph or \subparagraph
\usepackage{etoolbox}
\makeatletter
\patchcmd\longtable{\par}{\if@noskipsec\mbox{}\fi\par}{}{}
\makeatother
% Allow footnotes in longtable head/foot
\IfFileExists{footnotehyper.sty}{\usepackage{footnotehyper}}{\usepackage{footnote}}
\makesavenoteenv{longtable}
\setlength{\emergencystretch}{3em} % prevent overfull lines
\providecommand{\tightlist}{%
  \setlength{\itemsep}{0pt}\setlength{\parskip}{0pt}}
\usepackage{newunicodechar}
\newunicodechar{→}{\ensuremath{\rightarrow}}
\newunicodechar{⟨}{\ensuremath{\langle}}
\newunicodechar{⟩}{\ensuremath{\rangle}}
\newunicodechar{⇒}{\ensuremath{\Rightarrow}}
\newunicodechar{⟹}{\ensuremath{\Longrightarrow}}
\newunicodechar{×}{\ensuremath{\times}}
\newunicodechar{≈}{\ensuremath{\approx}}
\newunicodechar{∼}{\ensuremath{\sim}}
\newunicodechar{≤}{\ensuremath{\le}}
\newunicodechar{≥}{\ensuremath{\ge}}
\newunicodechar{≠}{\ensuremath{\ne}}
\newunicodechar{∫}{\ensuremath{\int}}
\newunicodechar{−}{\ensuremath{-}}
\newunicodechar{✗}{\ensuremath{\times}}
\newunicodechar{✓}{\ensuremath{\checkmark}}
\usepackage{bookmark}
\IfFileExists{xurl.sty}{\usepackage{xurl}}{} % add URL line breaks if available
\urlstyle{same}
\hypersetup{
  hidelinks,
  pdfcreator={LaTeX via pandoc}}

\author{}
\date{}

\begin{document}

\section{Emergence of Born Statistics from the Equidistribution of the
Observer--System Beat --- A Conjecture on the Localized-Kernel
Model}\label{emergence-of-born-statistics-from-the-equidistribution-of-the-observersystem-beat-a-conjecture-on-the-localized-kernel-model}

\textbf{Noriaki Kihara} (Derived from peer-review dialogue notes.
\textbf{This is neither a proof paper nor an observation paper, but a
conjecture (prediction) paper.} It is not a claim of established
physics, but the presentation of a candidate mechanism, made explicit,
together with its testable consequences.)

Version v0.4 (2026-06-28) DOI (Version): 10.5281/zenodo.20985736 (this
version, v0.4) / v0.3: 10.5281/zenodo.20981910 DOI (Concept):
10.5281/zenodo.20967081 Zenodo: https://zenodo.org/records/20985736

\textbf{Revision note (v0.4)}: Correction of the scope of the
prediction. Previously §7, the Abstract, and §9 described the
finite-\(N\) \(O(1/N)\) reproduction correction as a ``measurable
deviation from / falsifiable prediction of the Born rule,'' which
over-stated this paper's position. This paper does \textbf{not} claim
the beat mechanism to be the \textbf{root} of indeterminacy, and does
\textbf{not} positively assert that the Born rule breaks when the
relative-phase ratio is commensurate (non-irrational). Rather, it adopts
the hypothesis that a \textbf{more fundamental indeterminacy, not
identified in this paper, may prevent the very realization of an
observation at a strictly integer (fully coherent) ratio} --- such a
mechanism \textbf{may exist} (this is a hypothesis; its existence is not
asserted) --- so that, \textbf{if it does}, the Born rule of standard
physics holds \textbf{more strongly}. Hence the \textbf{non-observation}
of the \(O(1/N)\) correction is consistent with this conjecture and is
\textbf{not} a refutation. Open problem O5 (a more fundamental
indeterminacy) is added to §8. \textbf{The mechanism, equations, and
numerical results are unchanged} (only the register and scope of the
prediction are corrected).

\textbf{Revision note (v0.3)}: (1) Removed the §7 clause ``the gradient
energy \(\sum\nu^2\) becomes exactly finite,'' which suggested an energy
divergence (the falsifiable prediction itself --- the \(O(1/N)\)
reproduction correction --- is unchanged). (2) Made explicit at the end
of §2 that the \textbf{match of frequency is self-evident} (the
localized wave is sculpted in \textbf{shape} by odd harmonics but its
fundamental frequency \(\nu=1\) is unchanged = scale invariance),
sharpening the division of labor between Point 1 (self-evident frequency
match + non-trivial shape reproduction) and Point 2 (non-trivial
fluctuation = beat). (3) Echoed in the §9 conclusion the role separation
``\textbf{Weyl equidistribution supplies only a uniform measure; the
\(|\psi|^2\) bias comes from the reproducing kernel + the bridge}''
(previously only in §5(D)). Claims, conclusions, and equations are
unchanged.

\textbf{Revision note (v0.2)} (reflecting the second review, Claude
Code): (1) both input cases derived explicitly --- localized input and
spread input \(\cos\varphi\) come from the same reproducing kernel, the
output difference being only the input difference (§3, §5(B), confirmed
by actual convolution to machine precision); (2) §3(I) role separation
rewritten --- equidistribution (Weyl) supplies a uniform measure (no
bias), the \(|\psi|^2\) shape is the deterministic reproducing-kernel
read, intensity→click is the bridge (§3 II); the conflation ``Weyl alone
makes the squared bias'' is removed; (3) falsification battery (§5 E)
--- spread input not peaked, localized input peaked, Weyl alone flat,
both cases one identity --- all four pass; (4) §5(B) upgraded from a
restatement of {[}2{]} to a self-contained verification by actual
convolution (tautology removed), with the band condition
(\(N\)-dependence) shown; (5) the confirmed DOI of {[}2{]} added; (6)
the ``proven core'' made conditional on the equidistribution assumption
(§8 O1), and the discrete-clock dependency (B1 upstream of B2/B3) made
explicit (§2, §8 O2); (7) the mononym ``Lynnx'' in {[}7{]} confirmed as
listed on arXiv; (8) the status of the §5(D) bridge Monte-Carlo made
explicit --- it is a consistency demonstration that \emph{assumes} the
bridge (§3 II), not a derivation/verification of it. \textbf{Claims,
conclusions, and equations are unchanged} (rigor and honesty only).

\begin{center}\rule{0.5\linewidth}{0.5pt}\end{center}

\subsection{Abstract}\label{abstract}

The preceding localized-kernel paper {[}2{]} showed, from the
\textbf{reproducing-kernel property} of the constant-amplitude
odd-harmonic sum (isolated peak wave \(S_N\)) on the half-wavelength
phase interval \(\varphi\in[-\pi/2,\pi/2]\), that the \textbf{form} of
the observed distribution coincides exactly with the base wave's
\(|\psi_{\rm base}|^2\). It left, however, ``the existence of
randomness'' --- why a deterministically computed interference-intensity
profile localizes to a single point each trial and appears as the
frequency \(|\psi_{\rm base}|^2\) only upon repetition --- as an
unsolved postulate.

This paper presents a candidate mechanism for this \textbf{origin of
randomness} as a \textbf{conjecture}. The central observation is that
the scan variable \(\varphi_0\) of the localized-kernel paper is
\textbf{not a free variable but the phase of the observer}. The observer
is itself a physical object with a finite frequency \(\psi\), and there
is no absolute phase reference outside the system (the subjective-space
premise). Hence the measurement sequence scans the \textbf{relative
phase} (beat/heterodyne) \((\nu-\psi)t\) of the system frequency \(\nu\)
and the observer frequency \(\psi\). When the relative-phase ratio is
irrational, this relative phase \textbf{equidistributes} over the
interval (Weyl's equidistribution theorem).

The logic has two stages. \textbf{(I) Conditionally proven core}: by the
reproducing-kernel property {[}2{]}, the squared overlap at each
observer phase \(\varphi_0\) is exactly
\((\pi/2)^2|\psi_{\rm base}(\varphi_0)|^2\), and \textbf{if} the beat
phase equidistributes (irrational ratio, §8 O1), the long-time
\textbf{position-resolved intensity profile} coincides, after
normalization, with \(|\psi_{\rm base}|^2\). What is guaranteed here is
the \textbf{shape of the deterministically swept intensity profile}, not
a single-shot probability. This consequence follows, granting the
equidistribution assumption (O1), from the reproducing-kernel property
(proven) and Weyl's theorem (proven) with no further assumption.
\textbf{(II) Conjecture (bridge)}: we conjecture the transition by which
this intensity profile appears, via threshold detection, as the
\textbf{frequency of single clicks} \(\propto|\psi_{\rm base}|^2\). (II)
is the part this paper does \textbf{not} derive; why a single point
stands up each shot (discreteness) is left unsolved here.

This conjecture is not an isolated idea but sits at the
\textbf{intersection} of two established lineages: (a) Khrennikov's
prequantum classical statistical field theory (PCSFT) and the emergence
of the Born rule via threshold detection {[}3--6{]} (the physical
picture), and (b) the \textbf{method of arbitrary functions}
{[}7--10{]}, begun by Poincaré and applied to the quantum case by
Feintzeig et al.~(the mathematical mechanism = our equidistribution).
The contribution specific to this paper is to physically identify the
``fast equidistributing random variable'' that both lineages place
abstractly as the observer--system \textbf{heterodyne beat}, and to fix
the shape of the intensity profile \textbf{exactly} by the
reproducing-kernel property.

This paper does \textbf{not} claim the beat mechanism to be the root of
indeterminacy. It is only one facet of indeterminacy, and the paper does
\textbf{not} positively assert that the Born rule breaks when the
relative-phase ratio is non-irrational (commensurate). At finite \(N\)
(discrete lattice) an \(O(1/N)\) reproduction correction may formally
appear; but this paper instead adopts the hypothesis that \textbf{a more
fundamental indeterminacy, not identified here, may prevent the very
realization of an observation at a strictly integer (fully coherent)
ratio --- such a mechanism may exist (its existence is not asserted)}
--- so that, \textbf{if it does}, the Born rule of standard physics
holds \textbf{more strongly}. Hence the non-observation of this
correction is consistent with the conjecture, and
\textbf{non-observation is not a refutation} (indeed, were a strictly
integer ratio realized and the correction observed, that would be a
surprise contrary to this paper's expectation).

\begin{center}\rule{0.5\linewidth}{0.5pt}\end{center}

\subsection{1. Background and the place of this
paper}\label{background-and-the-place-of-this-paper}

The localized-kernel paper {[}2{]} derived the ``form'' of the Born rule
\(P\propto|\psi|^2\) squared structure from the reproducing-kernel
property, narrowing the remaining postulates to two: (a) the
\textbf{squaring rule} of amplitude → probability (exponent 2), (c) the
\textbf{existence of randomness}. A \textbf{measurement model} (b)
representing measurement by the convolution of the localized kernel also
remains a premise. What {[}2{]} rigorously established is: for a
band-limited base wave
\(\psi_{\rm base}=\sum_{|k|\le N}c_k e^{ik\varphi}\), the convolution
with the shifted localized kernel \(S_N\) satisfies

\[
(S_N*\psi_{\rm base})(\varphi_0)=\frac{\pi}{2}\,\psi_{\rm base}(\varphi_0)
\quad\Longrightarrow\quad
\big|(S_N*\psi_{\rm base})(\varphi_0)\big|^2=\frac{\pi^2}{4}\,|\psi_{\rm base}(\varphi_0)|^2
\]

exactly over the whole interval and for all \(N\) (a genuine modulus
\(Z\bar Z\) in the complex case).

This paper takes up (c). It is the deepest of the three postulates, the
measurement problem itself, unsolved even in standard theory. We do
\textbf{not claim to have solved it}; we formulate a candidate mechanism
as a \textbf{conjecture}, show that part of it (the reproduction of the
intensity profile) follows from the proven core, and explicitly carve
out the rest (discreteness, the probability interpretation) as unsolved.

\begin{center}\rule{0.5\linewidth}{0.5pt}\end{center}

\subsection{\texorpdfstring{2. Sharpening the problem: who decides
\(\varphi_0\)?}{2. Sharpening the problem: who decides \textbackslash varphi\_0?}}\label{sharpening-the-problem-who-decides-varphi_0}

In {[}2{]}, \(\varphi_0\) is plotted over the whole interval as a free
variable. But \textbf{one observation occurs at only one \(\varphi_0\)}.
The deterministic computation gives the function
\(|\psi_{\rm base}(\cdot)|^2\), yet one measurement returns a single
realized value. This transition ``\textbf{function → a single
realization per trial}'' is the entire remaining mystery. We decompose
it into three logically distinct bridges:

\begin{itemize}
\tightlist
\item
  \textbf{B1 (discreteness)}: why only one result stands up each time
  (the whole profile does not appear at once).
\item
  \textbf{B2 (which one)}: granting discreteness, why which
  \(\varphi_0\) is unpredictable from the state.
\item
  \textbf{B3 (ensemble = intensity)}: why the trial frequency matches
  the intensity profile.
\end{itemize}

The weight \(\propto|\text{overlap}|^2\) (a golden-rule-like coupling
rate) is the least non-trivial of the three. Positional fluctuation is
the manifestation of B1·B2, not a mystery of the \textbf{form} of
\(|\psi|^2\) --- the form is already settled in {[}2{]} as the
intensity. This paper's conjecture mainly answers B2 and B3, and leaves
B1 unsolved.

However, the three bridges are \textbf{not fully independent}. The
equidistribution mechanism below (§3--§4) becomes non-trivial only on a
\textbf{discrete sequence of measurement events} --- in continuous time
the sweep of \(\varphi_0=(\nu-\psi)t\) becomes uniform without
irrationality and the rational counter-example disappears --- so B1
(discreteness) is logically \textbf{upstream} of the B2·B3 mechanism.
That is, the equidistribution argument tacitly presupposes a discrete
clock. This dependency is made explicit in §8 O2.

Note that the appearance of the observation at the same
\textbf{frequency} as the baseline \(\nu=1\) probability wave is itself
self-evident: although the localized wave is sculpted in \textbf{shape}
by odd harmonics, the fundamental frequency \(\nu=1\) of the composite
wave is unchanged (since \(\gcd(1,3,\dots,N)=1\), the fundamental period
is set by the lowest mode \(\nu=1\); this is merely an observation under
the scale-invariant relative normalization, made explicit in the
odd-harmonic localization paper). Hence what this paper asks is not the
frequency match, but \textbf{(a) the mechanism by which the observed
waveform is faithfully reproduced as \(|\psi_{\rm base}|^2\)} (the
reproducing-kernel property, {[}2{]}; non-trivial because a general
kernel would distort it) and \textbf{(b) the mechanism by which
observation fluctuates across trials} (the beat). Point 1 (frequency
match = self-evident, scale invariance / shape reproduction =
reproducing kernel) and Point 2 (fluctuation = beat) belong to different
layers.

\begin{center}\rule{0.5\linewidth}{0.5pt}\end{center}

\subsection{\texorpdfstring{3. Conjecture: equidistribution of
\(\varphi_0\) by the observer--system
beat}{3. Conjecture: equidistribution of \textbackslash varphi\_0 by the observer--system beat}}\label{conjecture-equidistribution-of-varphi_0-by-the-observersystem-beat}

\textbf{Premise (the only physics added).} The observer is a physical
object with a finite frequency \(\psi\). A perfect absolute phase
reference (an absolute clock outside the system) does not exist --- this
is a necessity of the subjective-space premise shared by {[}2{]} and
this paper (only differences are physical; absolute phase is
unobservable), not a new assumption. If the observer were a perfect
reference, \(\varphi_0\) would be fixed and the fluctuation would
vanish, but the premise forbids such a device.

\textbf{Beat.} The product of the system \(\cos(\nu t)\) and the
observer \(\cos(\psi t)\) is

\[
\cos(\nu t)\cos(\psi t)=\tfrac12\big[\cos((\nu-\psi)t)+\cos((\nu+\psi)t)\big]
\]

so the difference frequency \(\nu-\psi\) gives a slow envelope (the
fringe period). Because the observer phase advances at each measurement
timing, the effective scan phase is \(\varphi_0=(\nu-\psi)t\). If the
observer has the same frequency as the system (\(\psi=\nu\)),
\(\varphi_0\) is fixed and the same phase is read each time --- no
fringes, a single value. If the observer differs (\(\psi\ne\nu\)),
\(\varphi_0\) is scanned and repetition sweeps out the distribution.

\textbf{Determinism of the single shot.} Importantly, the single
convolution readout at each \(\varphi_0\) is, by {[}2{]},
\textbf{deterministic}:
\((S_N*\psi_{\rm base})(\varphi_0)=\tfrac{\pi}{2}\psi_{\rm base}(\varphi_0)\)
is sharp and unambiguous. This is consistent with the fact that the
observed peak wave is actually sharply localized (it does not scatter
randomly). The ambiguity resides not in the single shot but in the
\textbf{distribution of \(\varphi_0\) over repetitions}.

\textbf{Conjecture (main).} When measurements occur at discrete times
\(t_n\) (by the observer clock) and the relative-phase ratio
\(r=(\nu-\psi)\,\Delta t/(2\pi)\) is irrational,
\(\varphi_0=(\nu-\psi)t_n\) \textbf{equidistributes} over
\([-\pi/2,\pi/2]\) (Weyl's equidistribution theorem for an irrational
rotation). The reliance on a discrete measurement sequence (in
continuous time irrationality is unnecessary and the counter-example
disappears) corresponds to the dependency of §2. Hence in the long run
each \(\varphi_0\) is uniformly sampled, and the long-time average of
the position-resolved squared readout is

\[
\boxed{\;\lim_{T\to\infty}\frac1T\int_0^T \big|(S_N*\psi_{\rm base})\big((\nu-\psi)t\big)\big|^2\,\delta_{\varphi_0}\,dt
\;\propto\; |\psi_{\rm base}(\varphi_0)|^2\;}
\]

i.e.~the \textbf{position-resolved intensity profile} coincides, after
normalization, with \(|\psi_{\rm base}|^2\). This equality follows,
\textbf{granting the equidistribution assumption (§8 O1), with no
further assumption}, from the reproducing-kernel property (the exact
result of §1) and Weyl equidistribution (a classical theorem) (confirmed
numerically in §5). What is guaranteed is the \textbf{shape} of the
intensity profile, not a single-shot probability (the bridge is below).

\textbf{Role separation (removing the conflation).} Here we rigorously
separate three contributions. \textbf{Equidistribution (Weyl) only
supplies the uniform measure \(d\varphi_0\); by itself it makes no
squared bias} --- if the observer phase \(\varphi_0\) is swept
uniformly, all positions are equally frequent (constant), which is not
the \(|\psi|^2\) bias. The \textbf{source of the per-position weight
(shape) \(|\psi_{\rm base}(\varphi_0)|^2\) is the deterministic
reproducing-kernel read at each \(\varphi_0\)} ({[}2{]}:
\((S_N*\psi_{\rm base})(\varphi_0)=\tfrac{\pi}{2}\psi_{\rm base}(\varphi_0)\),
reproducing the input \(\psi_{\rm base}\) as is), not equidistribution.
And what turns the intensity profile into the frequency of single clicks
is the \textbf{threshold-detection bridge} (§3 (II)). Schematically:

\[
\underbrace{\text{frequency of }|\psi_{\rm base}(\varphi_0)|^2}_{\text{Born output}}
=\underbrace{d\varphi_0\ (\text{uniform})}_{\text{Weyl equidistribution}}\ \times\
\underbrace{|\psi_{\rm base}(\varphi_0)|^2}_{\text{deterministic kernel read (shape)}}\ \times\
\underbrace{(\text{intensity}\to\text{click})}_{\text{bridge, §3 (II)}}.
\]

The \textbf{proven core} asserts only the middle factor (the
intensity-profile shape \(=|\psi_{\rm base}|^2\)), with Weyl supplying
the left uniform measure. The right bridge is conjecture. We do not
write as if Weyl alone makes the bias (in §5 D we also confirm
numerically that the histogram of equidistributed \(\varphi_0\) is
\textbf{flat}, and \(|\psi|^2\) appears only through intensity
weighting).

\textbf{Exact derivation of both cases (same mechanism, input-dependent
only).} Since the reproducing-kernel identity
\((S_N*\psi_{\rm base})(\varphi_0)=\tfrac{\pi}{2}\psi_{\rm base}(\varphi_0)\)
holds independently of the input \(\psi_{\rm base}\), for two
representative inputs:

\begin{itemize}
\tightlist
\item
  \textbf{Case 1 (localized)}: for
  \(\psi_{\rm base}^{\rm loc}=\sum_{|k|\le N}c_k e^{ik\varphi}\) (a
  multimode superposition localized at the center; the most localized
  example is \(S_N\) itself),
  \[\big|(S_N*\psi_{\rm base}^{\rm loc})(\varphi_0)\big|^2=\tfrac{\pi^2}{4}\,|\psi_{\rm base}^{\rm loc}(\varphi_0)|^2\ (\text{peaked output}).\]
\item
  \textbf{Case 2 (spread)}: for \(\psi_{\rm base}^{\cos}=\cos\varphi\)
  (a spread single mode),
  \[\big|(S_N*\cos)(\varphi_0)\big|^2=\tfrac{\pi^2}{4}\cos^2\varphi_0\ (\text{spread output}).\]
\end{itemize}

Both come from the \textbf{same reproducing kernel}, and the output
difference is only the input \(\psi_{\rm base}\) difference. \textbf{A
spread input (cos) gives a spread output} (\(\cos^2\)); \textbf{a
localized input gives a peaked output} --- since the reproducing kernel
reproduces the input without distortion, there is no breakdown beyond
the input dependence. Case 2 is numerically supported by Fig. 1(a) of
{[}2{]} (the cos base lies on \(\cos^2\) for all \(N\)), and is included
as an explicit test in this paper's verification code (§5). This
\textbf{agreement of both cases} is the strongest evidence that the
mechanism is consistent with input dependence only (clearing all
falsification conditions) (§5 E).

\textbf{Conjecture (bridge, unproven).} The above intensity profile
appears, via threshold detection, as the \textbf{frequency of single
clicks} \(\propto|\psi_{\rm base}|^2\). That is, B3 (frequency =
intensity) is reduced to the intensity profile, and B1 (single-shot) is
delegated to threshold detection. This bridge is the part this paper
does \textbf{not} derive, connecting to the mechanism carried by lineage
A (Khrennikov) of §6.

\begin{center}\rule{0.5\linewidth}{0.5pt}\end{center}

\subsection{4. Why determinism and probability coexist: the method of
arbitrary
functions}\label{why-determinism-and-probability-coexist-the-method-of-arbitrary-functions}

The mathematical body of the mechanism of §3 is the \textbf{method of
arbitrary functions}. This is the classical-probability tradition in
which a fast-varying phase (or dynamical parameter) and \textbf{any
smooth distribution} over it converge, in a certain limit, to a
\textbf{universal output distribution} independent of the input details
{[}8--10{]}. Poincaré analyzed the rotation of a roulette wheel and
showed that any smooth distribution of the initial angular velocity
gives nearly equal probability to each pocket at high speed {[}8{]}.
Diaconis et al.~gave the same kind of result for coin tossing {[}17{]}.
The point of the method is to \textbf{reconcile non-trivial
probabilities (neither 0 nor 1) with a deeper layer of determinism}.

Our setup is the quantum version of this. Just as fast rotation
equidistributes over the \textbf{alternating red-black arcs} of
Poincaré's wheel (a periodic binary structure on the circle), so our
measurement reads a \textbf{0/180 binary sign} (the sign of the
real-axis projection), and the high-frequency beat equidistributes the
relative phase --- the same mathematical object: equidistribution of a
fast phase over a periodic structure. Feintzeig et al.~{[}7{]} showed
that treating a dynamical parameter as a random variable with an initial
distribution yields the Born rule as \textbf{universal limiting
behavior}, for any initial distribution, in the joint long-time and
classical limit. This paper physically identifies that
``equidistributing dynamical parameter'' as the observer--system
\textbf{relative phase}.

Therefore randomness is intrinsic neither to the state nor to the
observer alone. It emerges from the \textbf{relative phase of two
deterministic oscillators} equidistributing over the measurement
sequence. This breaks the circularity: \(|\psi|^2\) is assumed nowhere;
the fluctuation is produced from the independent fact of the boundary of
observability (the absence of absolute phase).

\begin{center}\rule{0.5\linewidth}{0.5pt}\end{center}

\subsection{5. Numerical corroboration (the testable
core)}\label{numerical-corroboration-the-testable-core}

We confirm numerically the \textbf{(I) proven core} of §3. Base wave
\(\psi_{\rm base}=\cos\varphi+0.5\cos3\varphi-0.3\cos5\varphi\),
\(N=9\).

\textbf{(A) Weyl equidistribution of the beat phase.} Taking the
relative-phase ratio irrational (golden ratio \((\sqrt5-1)/2\)), the
star-discrepancy from the uniform distribution when \(\varphi_0\) is
sampled by the irrational rotation:

{\def\LTcaptype{none} % do not increment counter
\begin{longtable}[]{@{}
  >{\raggedright\arraybackslash}p{(\linewidth - 8\tabcolsep) * \real{0.2000}}
  >{\raggedright\arraybackslash}p{(\linewidth - 8\tabcolsep) * \real{0.2000}}
  >{\raggedright\arraybackslash}p{(\linewidth - 8\tabcolsep) * \real{0.2000}}
  >{\raggedright\arraybackslash}p{(\linewidth - 8\tabcolsep) * \real{0.2000}}
  >{\raggedright\arraybackslash}p{(\linewidth - 8\tabcolsep) * \real{0.2000}}@{}}
\toprule\noalign{}
\begin{minipage}[b]{\linewidth}\raggedright
beats \(N_{\rm beats}\)
\end{minipage} & \begin{minipage}[b]{\linewidth}\raggedright
\(10^2\)
\end{minipage} & \begin{minipage}[b]{\linewidth}\raggedright
\(10^3\)
\end{minipage} & \begin{minipage}[b]{\linewidth}\raggedright
\(10^4\)
\end{minipage} & \begin{minipage}[b]{\linewidth}\raggedright
\(10^5\)
\end{minipage} \\
\midrule\noalign{}
\endhead
\bottomrule\noalign{}
\endlastfoot
star-discrepancy & \(1.4\times10^{-2}\) & \(1.3\times10^{-3}\) &
\(2.6\times10^{-4}\) & \(3.0\times10^{-5}\) \\
\end{longtable}
}

Converges to \(0\) as \(N_{\rm beats}\) grows (equidistribution). By
contrast, for \textbf{rational ratios} (\(1/2,1/3,1/4\)) the discrepancy
is stuck at \(0.25\)--\(0.50\) and does not equidistribute (the fringe
is biased; counter-example). That is, ``the observer has a generic
(irrational-ratio) frequency relative to the system'' is the condition
for equidistribution.

\textbf{(B) Actual convolution of the reproducing kernel (both cases).}
In v0.2 we \textbf{actually integrate}
\((S_N*\psi_{\rm base})(\varphi_0)=\int_{-\pi/2}^{\pi/2}S_N(\varphi_0-\varphi)\psi_{\rm base}(\varphi)\,d\varphi\)
and confirm agreement with \(\tfrac{\pi}{2}\psi_{\rm base}(\varphi_0)\)
self-containedly in this paper's code (not merely restating {[}2{]}).
Machine-precision agreement for all three representative inputs:

{\def\LTcaptype{none} % do not increment counter
\begin{longtable}[]{@{}
  >{\raggedright\arraybackslash}p{(\linewidth - 4\tabcolsep) * \real{0.3333}}
  >{\raggedright\arraybackslash}p{(\linewidth - 4\tabcolsep) * \real{0.3333}}
  >{\raggedright\arraybackslash}p{(\linewidth - 4\tabcolsep) * \real{0.3333}}@{}}
\toprule\noalign{}
\begin{minipage}[b]{\linewidth}\raggedright
input \(\psi_{\rm base}\)
\end{minipage} & \begin{minipage}[b]{\linewidth}\raggedright
\(\max|(S_N*\psi)-\tfrac{\pi}{2}\psi|\)
\end{minipage} & \begin{minipage}[b]{\linewidth}\raggedright
\(\max\big||\,\cdot\,|^2_{\rm norm}-|\psi|^2_{\rm norm}\big|\)
\end{minipage} \\
\midrule\noalign{}
\endhead
\bottomrule\noalign{}
\endlastfoot
Case 1 localized \(\psi^{\rm loc}=S_N\) & \(1.8\times10^{-15}\) &
\(8.9\times10^{-16}\) \\
Case 1' multimode \(\cos+0.5\cos3-0.3\cos5\) & \(4.4\times10^{-16}\) &
\(4.4\times10^{-16}\) \\
Case 2 spread \(\cos\varphi\) & \(4.4\times10^{-16}\) &
\(3.3\times10^{-16}\) \\
\end{longtable}
}

From the \textbf{same reproducing kernel}, the output difference is only
the input difference: localized input \(\to\) peaked output, spread
input \(\to\) spread output (numerical confirmation of the both-case
derivation of §3).

\textbf{Band condition (\(N\)-dependence).} The reproducing kernel
\(S_N\) reproduces only modes \(|k|\le N\). Verifying by varying \(N\)
(each cell is \(\max|(S_N*\psi)-\tfrac{\pi}{2}\psi|\)):

{\def\LTcaptype{none} % do not increment counter
\begin{longtable}[]{@{}llllll@{}}
\toprule\noalign{}
input & \(N=1\) & \(N=3\) & \(N=5\) & \(N=9\) & \(N=31\) \\
\midrule\noalign{}
\endhead
\bottomrule\noalign{}
\endlastfoot
\(\cos\varphi\) (mode 1) & 2e-16 & 2e-16 & 4e-16 & 4e-16 & 2e-16 \\
multimode (1,3,5) & 1.2 & 0.47 & 4e-16 & 4e-16 & 7e-16 \\
\end{longtable}
}

\(\cos\) is \textbf{exact for all \(N\)}; the multimode (highest degree
5) is \textbf{exact for \(N\ge5\) and dropped for \(N<5\)}. This is the
manifestation of premise (ii) ``ignore unobservable phase differences''
= truncation exact once it covers the band. Note (B) is not the previous
``identity by substituting \((\pi/2)\psi\)'' but a
\textbf{self-contained independent verification by actually convolving
\(S_N\)} (the previous tautology is removed).

\textbf{(C) Single complex mode.} For \(\psi=e^{ik\varphi}\)
(\(k=1,3,5\)) the position-dependent spread of the squared readout is
\(\sim3\times10^{-15}\) (=0, uniform \(=|e^{ik\varphi}|^2=1\)). A sharp
single-frequency mode carries no position information and is fully
delocalized --- a manifestation of
\(\Delta\nu\cdot\Delta\varphi\gtrsim1\).

\textbf{(D) Numerical separation of roles (Weyl alone makes no bias).}
We directly check the role separation of §3. The \textbf{histogram
itself} of \(\varphi_0\) equidistributed by the irrational rotation
(\(2\times10^6\) points) is \textbf{flat}, with a maximum deviation from
the uniform density \(1/\pi\) of \(1.0\times10^{-5}\) ---
\textbf{equidistribution makes no bias}. On the other hand, when each
\(\varphi_0\) visit is accepted as a click with probability proportional
to the intensity \(|\psi_{\rm base}(\varphi_0)|^2\) (a Monte-Carlo model
of the threshold-detection bridge), the histogram of accepted clicks
matches \(|\psi_{\rm base}|^2\) (Case 2 cos: deviation
\(2.7\times10^{-3}\); Case 1' multimode: \(6.7\times10^{-3}\); both MC
statistical error). That is, the \(|\psi|^2\) bias comes from the
\textbf{deterministic kernel read + the bridge}, not from
equidistribution.

\textbf{Important (status statement).} This D2 Monte-Carlo is a
\textbf{consistency demonstration that assumes the bridge (§3 II) from
the outset} (``make clicks with probability proportional to the
intensity \(|\psi_{\rm base}(\varphi_0)|^2\)''); it is \textbf{neither a
derivation nor a verification of the bridge}. If you accept in
proportion to intensity, getting \(|\psi|^2\) out is automatic; what is
shown here is only the consistency that ``assuming the bridge closes the
whole thing without contradiction (uniform measure × intensity profile ×
bridge = \(|\psi|^2\) frequency).'' Hence the status separation of this
paper does not collapse: the \textbf{core} = Weyl (proven) × reproducing
kernel (proven) = uniform measure × the \textbf{shape}
\(|\psi_{\rm base}|^2\) of the intensity profile (backed by D1 and (B));
the \textbf{bridge} = intensity → single-click probability
(\textbf{conjecture}, in D2 merely assumed; delegated to threshold
detection = lineage A). D2 does \textbf{not} change the status that the
intensity is not yet a probability (the intensity profile is a classical
time-averaged interference fringe).

\textbf{(E) Falsification battery (consistency conditions).} We check
that the mechanism does not break with ``input dependence only,'' under
conditions that, if violated, would make the theory wrong:

\begin{enumerate}
\def\labelenumi{\arabic{enumi}.}
\tightlist
\item
  \textbf{Does a spread input give a peaked output?} The output
  peakedness (max/mean) of the cos input is \(1.96\) (low = stays
  spread). \textbf{OK} (no breakdown).
\item
  \textbf{Does a localized input over-spread?} The output peakedness of
  the \(S_N\) input is \(9.80\) (high = stays peaked). \textbf{OK}.
\item
  \textbf{Does Weyl alone make a squared bias?} Confirmed flat in (D1).
  \textbf{OK}.
\item
  \textbf{Can Case 2 be false independently of Case 1?} Both are special
  cases of the one identity \((S_N*\psi)=\tfrac{\pi}{2}\psi\) and cannot
  be refuted independently (on the same bridge). \textbf{OK}.
\end{enumerate}

All four conditions pass. This is a confirmation that ``both cases come
out correctly with the same mechanism and input dependence only,'' and
is \textbf{positive (offensive) evidence for the consistency of the
conjecture}.

The product of (A) (Weyl = uniform measure) and (B) (reproducing kernel
= deterministic read of the shape) constitutes the boxed equality of §3
--- the \textbf{shape} of the position-resolved intensity profile
\(=|\psi_{\rm base}|^2\). This is the paper's \textbf{(conditional on
the equidistribution assumption O1) proven core}. (D)(E) corroborate the
role separation and the consistency. The remaining bridge (intensity →
single-shot probability) is a claim of a kind not confirmable
numerically, treated in §6·§8.

\begin{center}\rule{0.5\linewidth}{0.5pt}\end{center}

\subsection{6. Relation to prior work}\label{relation-to-prior-work}

This conjecture sits at the intersection of two established lineages.

\subsubsection{6.1 Lineage A: Khrennikov's PCSFT + threshold detection
(the physical
picture)}\label{lineage-a-khrennikovs-pcsft-threshold-detection-the-physical-picture}

In Khrennikov's prequantum classical statistical field theory, a quantum
system is a symbolic representation of a classical random field;
Schrödinger dynamics is a special form of the linear dynamics of
classical fields; measurement is the interaction of the classical field
with a detector; and ``wave-packet collapse'' = discontinuity has the
\textbf{trivial origin of the use of a threshold detector} {[}6{]}. The
detection threshold is set by the average energy of the prequantum
signal, and the Born rule is reproduced as the measurement of a
classical signal by a properly calibrated threshold detector {[}3,4{]}.
The presence of a background field (vacuum fluctuation) is a key element
{[}4{]}.

Agreement with this paper: \textbf{discrete detection comes from the
detector, not the state} (B1); \textbf{the measuring device has its own
physics}. Difference: Khrennikov's source of randomness is a
\textbf{stochastic prequantum field} (\(L^2\)-valued Wiener process /
Gaussian field) {[}5{]}, and the Born rule is \textbf{approximate} and
predicts violations. Our source of randomness is the
\textbf{deterministic beat} (equidistribution of the relative phase),
and by the reproducing-kernel property the form is \textbf{exact}. The
``bridge'' of our conjecture (intensity → single click) corresponds
exactly to the mechanism carried by Khrennikov-type threshold detection.
That is, this paper cites lineage A as a \textbf{supplier of the
mechanism}.

\subsubsection{6.2 Lineage B: the method of arbitrary functions (the
mathematical
mechanism)}\label{lineage-b-the-method-of-arbitrary-functions-the-mathematical-mechanism}

As in §4, our equidistribution argument is the method of arbitrary
functions itself, begun by Poincaré {[}8{]}, Hopf {[}9{]}, Engel
{[}10{]}, with the quantum application by Feintzeig et al.~{[}7{]}. The
contribution specific to this paper is to physically identify the
``dynamical parameter taken as a random variable,'' which they place
abstractly, as the observer--system \textbf{heterodyne relative phase}
(extending the line by which {[}2{]} grounded it in the
Ramsey/heterodyne metrology of {[}15{]}), and to fix the shape of the
intensity profile exactly by the reproducing-kernel property.

\subsubsection{6.3 Adjacent lineages and
typicality}\label{adjacent-lineages-and-typicality}

On the foundational side, 't Hooft's deterministic (cellular-automaton)
quantum mechanics {[}12{]} stands alongside as a discrete substrate, and
Buniy--Hsu--Zee ``Discreteness and the origin of probability in quantum
mechanics'' {[}11{]} belongs to the discretization-first stance. Also
Zurek's envariance {[}13{]} and the typicality of Everett-type accounts
belong to the \textbf{same ``symmetry/typicality → weight'' family} as
the method of arbitrary functions. Our equidistribution (B3), the equal
amplitude (reproducing kernel) of {[}2{]}, and the equiprobable branches
of envariance can be seen as different faces of one family.

\subsubsection{6.4 Summary of placement}\label{summary-of-placement}

The contribution of this paper is clear as an intersection:
\textbf{Khrennikov-type measurement-induced discreteness (lineage A) +
the universal equidistribution weight of the method of arbitrary
functions (lineage B) + the exactness of the form by the reproducing
kernel ({[}2{]}) + the heterodyne identification of the equidistributing
fast variable (this paper)}. This combination is not found in the
existing literature.

\begin{center}\rule{0.5\linewidth}{0.5pt}\end{center}

\subsection{\texorpdfstring{7. The place of the finite-\(N\) correction:
the hypothesis that the Born rule holds more
strongly}{7. The place of the finite-N correction: the hypothesis that the Born rule holds more strongly}}\label{the-place-of-the-finite-n-correction-the-hypothesis-that-the-born-rule-holds-more-strongly}

It is suggestive that lineages A and B both yield the Born rule only in
a \textbf{limit/approximation}: Khrennikov approximately (and predicts
violations) {[}3{]}, the method of arbitrary functions in the joint
long-time and classical limit {[}7{]}. Our conjecture too is stated as
the universal/typical behavior of the Born rule in the
\textbf{equidistribution limit} (irrational beat, large \(N\)).
Formally, replacing the continuous-interval kernel \(S_N\) of {[}2{]} by
an \textbf{odd-harmonic kernel on a finite lattice} (a kernel hard-cut
at the Nyquist of discrete spacetime), the reproduction is no longer
exact and carries an \(O(1/N)\) correction. Here \(N=R/\ell\) is the
scale ratio (the system localization width \(\ell\) and the interval
\(R\)), and the size of the correction is \(\sim1/N=\ell/R\), largest at
the band edge (near the high-order modes).

\textbf{However, unlike lineage A (Khrennikov), this paper does not
present this correction as an observable deviation from the Born rule.}
The position of this paper is as follows. The beat mechanism is one
\textbf{facet} of indeterminacy, not its \textbf{root}. If the
relative-phase ratio became a strictly integer (fully coherent) one, the
equidistribution would formally break and the correction would surface.
But this paper adopts the hypothesis that \textbf{a more fundamental
indeterminacy, not identified here, may prevent the very realization of
such a strictly integer ratio --- such a mechanism may exist}; this is
only a hypothesis, and the paper does not assert the existence of such a
mechanism. \textbf{If it does exist}, the operation of fixing observer
and system to be fully coherent at arbitrary precision is in principle
unreachable, and as a consequence the Born rule of standard physics
holds, rather, \textbf{more strongly}.

Therefore this paper does \textbf{not} present the \(O(1/N)\) correction
as a ``measurable deviation'' or a ``falsifiable prediction.''
\textbf{The non-observation of the correction is consistent with this
conjecture, and non-observation is not a refutation} (indeed, were a
strictly integer ratio realized and the correction observed, that would
be a surprise contrary to this paper's expectation). What this paper
positively claims is not a breakdown of the Born rule, but that the beat
mechanism provides one facet of indeterminacy, and that the existence of
a more fundamental indeterminacy is left as an open problem (§8 O5).

\begin{center}\rule{0.5\linewidth}{0.5pt}\end{center}

\subsection{8. Open problems}\label{open-problems}

We make explicit what this paper does \textbf{not} derive.

\begin{itemize}
\tightlist
\item
  \textbf{O1 (the equidistribution assumption)}: for the
  position-resolved intensity to match \(|\psi|^2\), the relative phase
  must equidistribute with an irrational ratio (the condition confirmed
  in §5). This is natural but an assumption, and B3 is not fully
  eliminated but \textbf{relocated} to ``equidistribution of the
  observer phase.'' Still, this is far weaker than directly assuming
  \(|\psi|^2\).
\item
  \textbf{O2 (B1, single-shot)}: the beat gives the
  \textbf{distribution} of \(\varphi_0\) but does not give why a
  \textbf{single} click stands up per trial. This is a \textbf{separate
  mechanism} outside this paper's scope, also a part carried by lineage
  A (threshold detection). Note that the equidistribution argument of
  §3--§4 (irrational rotation → Weyl) presupposes a \textbf{discrete
  sequence of measurement events} (in continuous time the sweep is
  uniform without irrationality, and the rational counter-example
  disappears). \textbf{This paper leaves this discrete measurement
  sequence as an assumption and does not enter into its origin (the
  discretization mechanism).} Hence the discrete clock of B1 is also a
  \textbf{premise} of the B2·B3 mechanism, and O2 is not fully
  independent of the other open problems. Equidistribution (O1) and
  single-shot-ness (O2) share the common upstream of ``the existence of
  discrete measurement events.''
\item
  \textbf{O3 (the power of the square)}: the \textbf{exponent 2} that
  turns amplitude into probability should come from the area measure
  (Robertson-type) of the complex structure (real 2D), and does not come
  from the beat. It belongs to the complex-modulus argument
  \(Z\bar Z=\mathrm{Re}^2+\mathrm{Im}^2\) of {[}2{]}.
\item
  \textbf{O4 (scale attribution)}: in a map of type
  \(\sum\nu_n^2=\nu^2\), one must fix which of \(\nu\) and the internal
  \(\nu_n\) is the higher scale. The single frequency \(\nu\) seen by
  the observer and the direction of screening of the internal fine phase
  depend on this attribution.
\item
  \textbf{O5 (a more fundamental indeterminacy)}: the beat mechanism of
  this paper is one \textbf{facet} of indeterminacy, not its
  \textbf{root}. This paper does not positively assert that the Born
  rule breaks when the relative-phase ratio becomes a strictly integer
  (fully coherent) one; rather, it adopts the hypothesis that \textbf{a
  more fundamental indeterminacy may exist that prevents the very
  realization of such a fully coherent observation and screens the
  \(O(1/N)\) correction} --- its existence is not asserted. This paper
  cannot identify this fundamental indeterminacy and leaves it unsolved.
  The basis of the §7 hypothesis that the Born rule of standard physics
  holds \textbf{more strongly} also belongs here.
\end{itemize}

\begin{center}\rule{0.5\linewidth}{0.5pt}\end{center}

\subsection{9. Conclusion}\label{conclusion}

This paper presented, as a \textbf{conjecture}, a candidate mechanism
for the \textbf{origin of randomness} left as a postulate by the
localized-kernel paper {[}2{]}. The central observation is that the scan
phase \(\varphi_0\) is the observer's phase; that, since the observer is
a physical object with a finite frequency, no absolute reference exists
(the subjective-space premise); hence the measurement sequence scans the
observer--system relative phase (the beat), which equidistributes with
an irrational ratio (Weyl).

The logic is two-stage. \textbf{(Conditional on O1) proven core}: from
the reproducing-kernel property ({[}2{]}) and Weyl equidistribution, the
\textbf{position-resolved intensity profile} swept by the beat coincides
with \(|\psi_{\rm base}|^2\) (granting equidistribution O1, with no
further assumption; confirmed numerically in §5). \textbf{Conjectured
bridge}: this intensity appears, via threshold detection, as the
frequency of single clicks \(\propto|\psi_{\rm base}|^2\) (this paper
does not derive it; it connects to lineage A). Here Weyl
equidistribution supplies only a \textbf{uniform measure (the locus of
the fluctuation)}; the \(|\psi_{\rm base}|^2\) bias comes from the
\textbf{deterministic reproducing-kernel read + the bridge} (§5 D). That
the observation coincides with the baseline frequency \(\nu=1\) is
itself self-evident, since the harmonics change only the waveform and
not the fundamental frequency (scale invariance); what is non-trivial is
the \textbf{faithful reproduction of the waveform (reproducing kernel)
and the origin of the fluctuation (beat)}.

This conjecture sits at the intersection of two established lineages ---
Khrennikov's PCSFT + threshold detection (the physical picture) and the
method of arbitrary functions (the mathematical mechanism) --- and its
specific contribution is to physically identify the equidistributing
fast variable as the observer--system heterodyne beat and to make the
form exact by the reproducing-kernel property. At finite \(N\) (discrete
lattice) an \(O(1/N)\sim\ell/R\) reproduction correction may formally
appear, but this paper does not claim it as a measurable deviation. The
beat mechanism is one facet of indeterminacy, not its root; this paper
adopts the hypothesis that a more fundamental indeterminacy, not
identified here, may prevent the realization of an observation at a
strictly integer (fully coherent) ratio --- such a mechanism may exist
(its existence is not asserted) --- so that, if it does, the Born rule
of standard physics holds more strongly. Hence the non-observation of
this correction is also consistent with the conjecture (§7, §8 O5).

What remains is B1 (single-shot), the exponent 2, the scale attribution,
and a more fundamental indeterminacy (§8 O5), which this paper honestly
carves out as unsolved. We have neither ``derived'' the Born rule nor
``predicted its violation''; the scope of this paper is to present a
\textbf{candidate that relocates one facet of its statistical origin to
two independently correct ingredients --- the finiteness of the observer
and the equidistribution of the relative phase}.

\begin{center}\rule{0.5\linewidth}{0.5pt}\end{center}

\subsection{Appendix A: Sketch of the minimal
verification}\label{appendix-a-sketch-of-the-minimal-verification}

Base wave
\(\psi_{\rm base}=\cos\varphi+0.5\cos3\varphi-0.3\cos5\varphi\),
\(N=9\).

\begin{enumerate}
\def\labelenumi{\arabic{enumi}.}
\tightlist
\item
  \textbf{Weyl equidistribution}: generate
  \(\varphi_0^{(n)}=(r\,n \bmod 1)\cdot\pi-\pi/2\) with the
  relative-phase ratio \(r=(\sqrt5-1)/2\) (irrational) and confirm the
  star-discrepancy converges to \(0\) as \(N_{\rm beats}\to\infty\) (§5
  table). For rational \(r\) it does not converge (counter-example).
\item
  \textbf{Position-resolved intensity}: normalize
  \(|(\tfrac{\pi}{2})\psi_{\rm base}(\varphi_0)|^2\) at each
  \(\varphi_0\) and confirm machine-precision agreement with
  \(|\psi_{\rm base}|^2\).
\item
  \textbf{Single complex mode}: confirm that the squared readout is
  position-uniform (\(=1\)) for \(\psi=e^{ik\varphi}\).
\end{enumerate}

These are reproduced by the verification code as machine precision /
classical theorems. The bridge (intensity → single-shot probability) is
outside numerical verification and can be questioned only conceptually
(on the observability of the finite-\(N\) correction there is the
reservation of §7).

\begin{center}\rule{0.5\linewidth}{0.5pt}\end{center}

\subsection{References}\label{references}

{[}1{]} M. Born, ``Zur Quantenmechanik der Stoßvorgänge,'' \emph{Z.
Phys.} \textbf{37}, 863 (1926). {[}2{]} N. Kihara, ``Deriving the Form
of the Born Distribution from the Reproducing-Kernel Property of a
Localized Odd-Harmonic Wave,'' Zenodo, v0.1 (2026-06-27), Concept DOI:
10.5281/zenodo.20965526 / Version DOI: 10.5281/zenodo.20965527 (the
starting point: exact derivation of the form by the reproducing kernel).
{[}3{]} A. Khrennikov, ``Born's rule from measurements of classical
signals by threshold detectors which are properly calibrated,''
arXiv:1105.4269 (2011). {[}4{]} A. Khrennikov, ``Quantum Probabilities
and Violation of CHSH-Inequality from Classical Random Signals and
Threshold Type Detection Scheme,'' \emph{Prog. Theor. Phys.}
\textbf{128}, 31 (2012). {[}5{]} A. Khrennikov, ``Born's rule from
statistical mechanics of classical fields: from hitting times to quantum
probabilities,'' arXiv:1212.0756 (2012). {[}6{]} A. Khrennikov,
``Measurement problem: from de Broglie to theory of classical random
fields interacting with threshold detectors,'' arXiv:1301.0011 (2013).
{[}7{]} L. Bonds, B. Burson, K. Cicchella, B. H. Feintzeig, Lynnx, A.
Yusaini, ``Quantum Probability via the Method of Arbitrary Functions,''
arXiv:2409.16457 (2024). (The author ``Lynnx'' is a mononym as
registered on arXiv.) {[}8{]} H. Poincaré, \emph{Calcul des
probabilités} (Gauthier-Villars, Paris, 1896) (méthode des fonctions
arbitraires: the roulette wheel). {[}9{]} E. Hopf, ``On causality,
statistics and probability,'' \emph{J. Math. Phys.} \textbf{13}, 51
(1934). {[}10{]} E. Engel, \emph{A Road to Randomness in Physical
Systems}, Lecture Notes in Statistics \textbf{71} (Springer, 1992).
{[}11{]} R. V. Buniy, S. D. Hsu, A. Zee, ``Discreteness and the origin
of probability in quantum mechanics,'' \emph{Phys. Lett. B}
\textbf{640}, 219 (2006). {[}12{]} G. 't Hooft, \emph{The Cellular
Automaton Interpretation of Quantum Mechanics}, Fundamental Theories of
Physics \textbf{185} (Springer, 2016). {[}13{]} W. H. Zurek,
``Probabilities from entanglement, Born's rule \(p_k=|\psi_k|^2\) from
envariance,'' \emph{Phys. Rev.~A} \textbf{71}, 052105 (2005). {[}14{]}
A. Valentini and H. Westman, ``Dynamical origin of quantum
probabilities,'' \emph{Proc. R. Soc. A} \textbf{461}, 253 (2005).
{[}15{]} N. F. Ramsey, ``A Molecular Beam Resonance Method with
Separated Oscillating Fields,'' \emph{Phys. Rev.} \textbf{78}, 695
(1950). {[}16{]} H. Weyl, ``Über die Gleichverteilung von Zahlen mod.
Eins,'' \emph{Math. Ann.} \textbf{77}, 313 (1916) (the equidistribution
theorem). {[}17{]} P. Diaconis, S. Holmes, R. Montgomery, ``Dynamical
bias in the coin toss,'' \emph{SIAM Rev.} \textbf{49}, 211 (2007).

\begin{center}\rule{0.5\linewidth}{0.5pt}\end{center}

\emph{(This is a conjecture paper derived from peer-review dialogue
notes (Noriaki Kihara × Iris, 2026-06-27). It consists of a (conditional
on O1) proven core (the reproducing-kernel property + the
position-resolved intensity \(=|\psi_{\rm base}|^2\) via Weyl
equidistribution) and an explicitly demarcated unproven bridge
(intensity → single-click probability), the latter connecting to lineage
A (Khrennikov-type threshold detection). B1, the exponent 2, the scale
attribution, and a more fundamental indeterminacy (§8 O5) are left
unsolved.)}

\end{document}
